\section{Results}
\label{sec:results}

\subsection{Repeated cross-validation}
Table~\ref{tab:cv} and Fig.~\ref{fig:cv} summarise the 5$\times$10-fold comparison on the development set; Supplementary Table~S1 lists all nine metrics with confidence intervals for all 23 configurations. All learners except LightGBM and the single decision tree reached a mean ROC-AUC between 0.820 and 0.846, with overlapping corrected 95\% confidence intervals. Stacking on all 16 features had the highest mean ROC-AUC (0.846, 95\% CI 0.812--0.880), followed by stacking with the common RFECV subset (0.843, 0.809--0.878) and logistic regression on all features (0.842, 0.809--0.875). The differences between the stacks and logistic regression were small and not statistically significant (+0.004 and +0.001; Holm-adjusted $p = 1.000$ for both). After Holm adjustment, only LightGBM (0.803 and 0.800 without and with RFECV; adjusted $p = 0.044$ and $0.020$) and the decision tree (0.770 and 0.772; $p = 0.025$ and $0.020$) discriminated significantly worse than logistic regression. No configuration differed significantly from logistic regression in MCC.

At the default 0.5 threshold, the stacks had the highest mean sensitivity among the strong learners (0.763 without and 0.753 with RFECV, compared with 0.744 for logistic regression and 0.655 for the SVM), at the cost of lower specificity (0.757 compared with 0.775 and 0.840). The Platt-calibrated SVM had the lowest mean Brier score (0.156; stacking + RFECV 0.165).

The common RFECV subset changed mean ROC-AUC by at most 0.003 for every learner while reducing the inputs from 16 to a mean of 12.0 features (25\% fewer). Logistic regression with its own coefficient-ranked RFECV kept 10.1 features on average and reached a ROC-AUC of 0.840 (0.808--0.873). It performed no better than logistic regression with the common, random-forest-ranked subset (0.842) or with all features.

\begin{table*}[!htbp]
\centering
\caption{Development-set performance from 5$\times$10-fold repeated cross-validation (50 folds; $n = 614$). Values are means; Nadeau--Bengio corrected 95\% confidence intervals are shown in grey beneath. $\Delta$AUC is the mean paired difference from logistic regression on all features, with Holm-adjusted corrected resampled $t$-test $p$-values in brackets. ``+ RFECV'' denotes the common, random-forest-ranked RFECV subset. Sensitivity and specificity use a 0.5 threshold.}
\label{tab:cv}
\footnotesize
\setlength{\tabcolsep}{4pt}
\resizebox{\textwidth}{!}{%
\begin{tabular}{lrcccccl}
\toprule
Model & Features & ROC-AUC & MCC & Sens. & Spec. & Brier & $\Delta$AUC (adj.\ $p$) \\
\midrule
Stacking & 16.0 & 0.846 & 0.505 & 0.763 & 0.757 & 0.164 & +0.004 (1.000) \\
 &  & {\scriptsize\color{gray}(0.812--0.880)} & {\scriptsize\color{gray}(0.431--0.580)} &  &  &  &  \\[2pt]
Stacking + RFECV & 12.0 & 0.843 & 0.497 & 0.753 & 0.757 & 0.165 & +0.001 (1.000) \\
 &  & {\scriptsize\color{gray}(0.809--0.878)} & {\scriptsize\color{gray}(0.418--0.575)} &  &  &  &  \\[2pt]
Logistic regression & 16.0 & 0.842 & 0.507 & 0.744 & 0.775 & 0.163 & reference \\
 &  & {\scriptsize\color{gray}(0.809--0.875)} & {\scriptsize\color{gray}(0.438--0.576)} &  &  &  &  \\[2pt]
Logistic regression + RFECV & 12.0 & 0.842 & 0.492 & 0.717 & 0.782 & 0.163 & -0.000 (1.000) \\
 &  & {\scriptsize\color{gray}(0.811--0.873)} & {\scriptsize\color{gray}(0.418--0.566)} &  &  &  &  \\[2pt]
LR (LR-ranked RFECV) & 10.1 & 0.840 & 0.496 & 0.731 & 0.776 & 0.164 & -0.002 (1.000) \\
 &  & {\scriptsize\color{gray}(0.808--0.873)} & {\scriptsize\color{gray}(0.426--0.567)} &  &  &  &  \\[2pt]
SVM & 16.0 & 0.840 & 0.505 & 0.655 & 0.840 & 0.156 & -0.002 (1.000) \\
 &  & {\scriptsize\color{gray}(0.802--0.878)} & {\scriptsize\color{gray}(0.418--0.591)} &  &  &  &  \\[2pt]
SVM + RFECV & 12.0 & 0.837 & 0.483 & 0.643 & 0.830 & 0.158 & -0.006 (1.000) \\
 &  & {\scriptsize\color{gray}(0.800--0.873)} & {\scriptsize\color{gray}(0.393--0.572)} &  &  &  &  \\[2pt]
Extra trees + RFECV & 12.0 & 0.835 & 0.474 & 0.693 & 0.786 & 0.161 & -0.007 (1.000) \\
 &  & {\scriptsize\color{gray}(0.798--0.872)} & {\scriptsize\color{gray}(0.393--0.556)} &  &  &  &  \\[2pt]
Extra trees & 16.0 & 0.835 & 0.469 & 0.676 & 0.794 & 0.160 & -0.008 (1.000) \\
 &  & {\scriptsize\color{gray}(0.795--0.874)} & {\scriptsize\color{gray}(0.377--0.560)} &  &  &  &  \\[2pt]
Naive Bayes + RFECV & 12.0 & 0.832 & 0.473 & 0.624 & 0.836 & 0.205 & -0.010 (1.000) \\
 &  & {\scriptsize\color{gray}(0.797--0.867)} & {\scriptsize\color{gray}(0.393--0.554)} &  &  &  &  \\[2pt]
Random forest & 16.0 & 0.831 & 0.484 & 0.712 & 0.780 & 0.164 & -0.011 (1.000) \\
 &  & {\scriptsize\color{gray}(0.795--0.868)} & {\scriptsize\color{gray}(0.390--0.578)} &  &  &  &  \\[2pt]
Random forest + RFECV & 12.0 & 0.831 & 0.492 & 0.718 & 0.781 & 0.164 & -0.011 (1.000) \\
 &  & {\scriptsize\color{gray}(0.797--0.866)} & {\scriptsize\color{gray}(0.403--0.581)} &  &  &  &  \\[2pt]
Gradient boosting & 16.0 & 0.831 & 0.436 & 0.595 & 0.828 & 0.164 & -0.011 (1.000) \\
 &  & {\scriptsize\color{gray}(0.798--0.863)} & {\scriptsize\color{gray}(0.353--0.518)} &  &  &  &  \\[2pt]
Naive Bayes & 16.0 & 0.830 & 0.494 & 0.626 & 0.852 & 0.205 & -0.012 (1.000) \\
 &  & {\scriptsize\color{gray}(0.793--0.867)} & {\scriptsize\color{gray}(0.420--0.568)} &  &  &  &  \\[2pt]
Gradient boosting + RFECV & 12.0 & 0.830 & 0.441 & 0.601 & 0.827 & 0.164 & -0.012 (1.000) \\
 &  & {\scriptsize\color{gray}(0.798--0.862)} & {\scriptsize\color{gray}(0.354--0.528)} &  &  &  &  \\[2pt]
KNN & 16.0 & 0.824 & 0.466 & 0.611 & 0.841 & 0.162 & -0.019 (1.000) \\
 &  & {\scriptsize\color{gray}(0.781--0.866)} & {\scriptsize\color{gray}(0.375--0.558)} &  &  &  &  \\[2pt]
XGBoost & 16.0 & 0.821 & 0.436 & 0.611 & 0.816 & 0.173 & -0.021 (0.954) \\
 &  & {\scriptsize\color{gray}(0.784--0.857)} & {\scriptsize\color{gray}(0.346--0.527)} &  &  &  &  \\[2pt]
KNN + RFECV & 12.0 & 0.821 & 0.458 & 0.604 & 0.839 & 0.163 & -0.022 (1.000) \\
 &  & {\scriptsize\color{gray}(0.777--0.865)} & {\scriptsize\color{gray}(0.361--0.554)} &  &  &  &  \\[2pt]
XGBoost + RFECV & 12.0 & 0.820 & 0.440 & 0.615 & 0.817 & 0.172 & -0.023 (0.586) \\
 &  & {\scriptsize\color{gray}(0.785--0.854)} & {\scriptsize\color{gray}(0.357--0.524)} &  &  &  &  \\[2pt]
LightGBM & 16.0 & 0.803 & 0.414 & 0.622 & 0.789 & 0.205 & -0.039 (0.044) \\
 &  & {\scriptsize\color{gray}(0.764--0.842)} & {\scriptsize\color{gray}(0.321--0.508)} &  &  &  &  \\[2pt]
LightGBM + RFECV & 12.0 & 0.800 & 0.420 & 0.627 & 0.790 & 0.202 & -0.042 (0.020) \\
 &  & {\scriptsize\color{gray}(0.765--0.836)} & {\scriptsize\color{gray}(0.336--0.505)} &  &  &  &  \\[2pt]
Decision tree + RFECV & 12.0 & 0.772 & 0.429 & 0.767 & 0.679 & 0.207 & -0.070 (0.020) \\
 &  & {\scriptsize\color{gray}(0.731--0.813)} & {\scriptsize\color{gray}(0.348--0.510)} &  &  &  &  \\[2pt]
Decision tree & 16.0 & 0.770 & 0.431 & 0.774 & 0.675 & 0.210 & -0.072 (0.025) \\
 &  & {\scriptsize\color{gray}(0.729--0.811)} & {\scriptsize\color{gray}(0.352--0.511)} &  &  &  &  \\[2pt]
\bottomrule

\end{tabular}}
\end{table*}

\begin{figure}[!htbp]
  \centering
  \includegraphics[width=0.75\linewidth]{fig2_cv_auc.pdf}
  \caption{Cross-validated ROC-AUC of each learner with all 16 features (open circles) and with the common RFECV subset (filled circles). Bars are Nadeau--Bengio corrected 95\% confidence intervals over 50 folds. The dashed line marks logistic regression on all features. The stacking ensemble is highlighted.}
  \label{fig:cv}
\end{figure}

\FloatBarrier
\subsection{Selected features and their stability}
Across the 50 outer folds, random-forest-ranked RFECV kept a median of 14 features (interquartile range 8--15; range 7--16). Seven features were kept in every fold: glucose, BMI and the engineered \texttt{Glucose\_BMI}, \texttt{Glucose\_Age}, \texttt{BMI\_Age}, \texttt{HOMA\_IR\_surrogate} and \texttt{Pedigree\_Age} (Fig.~\ref{fig:selection}). The diabetes pedigree function was kept in 94\% of folds. The thresholded flags were the weakest candidates: \texttt{Obese} was kept in 6\% of folds and \texttt{IGT\_2h} in 38\%. They largely duplicate information already carried by BMI and glucose.

Coefficient-ranked RFECV for logistic regression preferred a different, smaller subset (median 10.5 features; range 3--14). It kept glucose in every fold, and age, pregnancies and the pedigree function in more than 90\% of folds. It rarely kept blood pressure (2\%) and never kept skin-fold thickness. A core of glycaemic, adiposity and diabetes-pedigree predictors was therefore retained by both selectors, but the exact subset depended on the ranking criterion.

\begin{figure}[!htbp]
  \centering
  \includegraphics[width=0.75\linewidth]{fig3_selection.pdf}
  \caption{Proportion of the 50 outer cross-validation folds in which each candidate feature was retained by random-forest-ranked RFECV (the common protocol) and by coefficient-ranked RFECV for logistic regression. Engineered features are shown in italics.}
  \label{fig:selection}
\end{figure}

\FloatBarrier
\subsection{Internal hold-out performance}
On the 154 hold-out patients (Table~\ref{tab:holdout}, Fig.~\ref{fig:rocpr}), the stacking + RFECV pipeline reached a ROC-AUC of 0.822 (95\% CI 0.756--0.888), sensitivity 0.815 (0.709--0.915), specificity 0.740 (0.649--0.823), MCC 0.532 (0.403--0.658) and F1 0.710. It correctly identified 44 of 54 patients with diabetes. Among the evaluated configurations it had the highest observed sensitivity, F1 and MCC on this split. However, the confidence intervals of all models overlapped widely. No ROC-AUC difference between the stacking + RFECV pipeline and any comparator was statistically significant (DeLong test; differences from $-0.004$ to $+0.013$; smallest Holm-adjusted $p = 0.89$). Supplementary Table~S2 gives all hold-out metrics. Random forest + RFECV (0.826) and XGBoost + RFECV (0.824) had slightly higher hold-out ROC-AUC, and average precision ranged from 0.683 (logistic regression) to 0.711 (random forest + RFECV).

\begin{table*}[!htbp]
\centering
\caption{Hold-out performance ($n = 154$, 54 with diabetes) at a 0.5 threshold. 95\% confidence intervals are shown in grey beneath: DeLong for ROC-AUC, percentile bootstrap (2000 resamples) otherwise. AP, average precision. The last column gives the ROC-AUC difference of stacking + RFECV minus the comparator, with Holm-adjusted DeLong $p$-values.}
\label{tab:holdout}
\footnotesize
\setlength{\tabcolsep}{3.5pt}
\resizebox{\textwidth}{!}{%
\begin{tabular}{lccccccl}
\toprule
Model & ROC-AUC & Sensitivity & Specificity & MCC & F1 & AP & $\Delta$AUC (adj.\ $p$) \\
\midrule
Stacking + RFECV & 0.822 & 0.815 & 0.740 & 0.532 & 0.710 & 0.696 & -- \\
 & {\scriptsize\color{gray}(0.756--0.888)} & {\scriptsize\color{gray}(0.709--0.915)} & {\scriptsize\color{gray}(0.649--0.823)} & {\scriptsize\color{gray}(0.403--0.658)} & {\scriptsize\color{gray}(0.615--0.797)} &  &  \\[2pt]
Stacking & 0.819 & 0.796 & 0.720 & 0.494 & 0.688 & 0.702 & +0.003 (1.000) \\
 & {\scriptsize\color{gray}(0.752--0.885)} & {\scriptsize\color{gray}(0.688--0.904)} & {\scriptsize\color{gray}(0.631--0.804)} & {\scriptsize\color{gray}(0.362--0.625)} & {\scriptsize\color{gray}(0.592--0.778)} &  &  \\[2pt]
Logistic regression & 0.811 & 0.741 & 0.740 & 0.464 & 0.667 & 0.683 & +0.011 (1.000) \\
 & {\scriptsize\color{gray}(0.744--0.878)} & {\scriptsize\color{gray}(0.625--0.862)} & {\scriptsize\color{gray}(0.649--0.825)} & {\scriptsize\color{gray}(0.325--0.606)} & {\scriptsize\color{gray}(0.564--0.763)} &  &  \\[2pt]
LR + LR-ranked RFECV & 0.809 & 0.722 & 0.740 & 0.447 & 0.655 & 0.684 & +0.013 (0.892) \\
 & {\scriptsize\color{gray}(0.741--0.876)} & {\scriptsize\color{gray}(0.608--0.846)} & {\scriptsize\color{gray}(0.649--0.825)} & {\scriptsize\color{gray}(0.309--0.591)} & {\scriptsize\color{gray}(0.551--0.753)} &  &  \\[2pt]
Random forest + RFECV & 0.826 & 0.685 & 0.760 & 0.434 & 0.643 & 0.711 & -0.004 (1.000) \\
 & {\scriptsize\color{gray}(0.761--0.892)} & {\scriptsize\color{gray}(0.564--0.804)} & {\scriptsize\color{gray}(0.673--0.840)} & {\scriptsize\color{gray}(0.290--0.573)} & {\scriptsize\color{gray}(0.537--0.739)} &  &  \\[2pt]
XGBoost + RFECV & 0.824 & 0.648 & 0.820 & 0.470 & 0.654 & 0.697 & -0.003 (1.000) \\
 & {\scriptsize\color{gray}(0.758--0.891)} & {\scriptsize\color{gray}(0.523--0.776)} & {\scriptsize\color{gray}(0.743--0.888)} & {\scriptsize\color{gray}(0.321--0.605)} & {\scriptsize\color{gray}(0.543--0.752)} &  &  \\[2pt]
SVM & 0.810 & 0.537 & 0.830 & 0.383 & 0.580 & 0.684 & +0.012 (1.000) \\
 & {\scriptsize\color{gray}(0.740--0.880)} & {\scriptsize\color{gray}(0.407--0.672)} & {\scriptsize\color{gray}(0.755--0.896)} & {\scriptsize\color{gray}(0.226--0.534)} & {\scriptsize\color{gray}(0.457--0.692)} &  &  \\[2pt]
\bottomrule

\end{tabular}}
\end{table*}

\begin{figure*}[!htbp]
  \centering
  \includegraphics[width=0.9\textwidth]{fig4_roc_pr.pdf}
  \caption{Hold-out ROC (a) and precision--recall (b) curves for stacking + RFECV, logistic regression and random forest + RFECV. The dotted lines show chance performance and the outcome prevalence (0.35).}
  \label{fig:rocpr}
\end{figure*}

\FloatBarrier
\subsection{Calibration}
All class-weighted models overestimated risk (Table~\ref{tab:calibration}, Fig.~\ref{fig:calibration}). On the pooled development-set out-of-fold predictions, the calibration intercept was $-0.616$ for stacking + RFECV and $-0.623$ for logistic regression. This is almost exactly the logit of the outcome prevalence, $\operatorname{logit}(0.349) = -0.62$, which is consistent with the expected effect of balanced class weighting: weighting the classes equally corresponds to training at a prevalence of 50\%, which shifts the model intercept by about $-\operatorname{logit}(\pi)$ for prevalence $\pi$. We did not run a controlled experiment without class weights, so this explanation is supported by the arithmetic and by the well-calibrated unweighted Platt layer of the SVM, but it is not formally isolated. The calibration slope of the stack was close to one (1.002), so its spread of risks was appropriate and only its level was shifted. The Platt-calibrated SVM, whose probabilities are refitted without class weights, was well calibrated (intercept $-0.029$, slope 0.985) and had the lowest Brier score. On the hold-out set, the stack's intercept ($-0.624$) was unchanged and its slope fell to 0.866, compared with 0.756 for logistic regression.

\begin{table}[!htbp]
\centering
\caption{Calibration. Development values use pooled out-of-fold predictions ($614 \times 5$); hold-out Brier scores have bootstrap 95\% CIs. Ideal values: intercept 0, slope 1.}
\label{tab:calibration}
\footnotesize
\setlength{\tabcolsep}{3pt}
\begin{tabular}{lcccccc}
\toprule
& \multicolumn{3}{c}{Development (out-of-fold)} & \multicolumn{3}{c}{Hold-out} \\
\cmidrule(lr){2-4}\cmidrule(lr){5-7}
Model & Brier & Intercept & Slope & Brier (95\% CI) & Intercept & Slope \\
\midrule
Stacking + RFECV & 0.165 & -0.616 & 1.002 & 0.177 {\scriptsize\color{gray}(0.143--0.213)} & -0.624 & 0.866 \\
Logistic regression & 0.163 & -0.623 & 0.899 & 0.181 {\scriptsize\color{gray}(0.147--0.219)} & -0.629 & 0.756 \\
Random forest + RFECV & 0.164 & -0.406 & 0.823 & 0.166 {\scriptsize\color{gray}(0.134--0.201)} & -0.338 & 0.776 \\
SVM & 0.156 & -0.029 & 0.985 & 0.169 {\scriptsize\color{gray}(0.135--0.205)} & +0.004 & 0.844 \\
\bottomrule

\end{tabular}
\end{table}

\begin{figure*}[!htbp]
  \centering
  \includegraphics[width=0.9\textwidth]{fig5_calibration.pdf}
  \caption{Calibration curves (deciles of predicted risk) on (a) pooled development-set out-of-fold predictions and (b) the hold-out set (quintiles). Points below the diagonal indicate overestimated risk.}
  \label{fig:calibration}
\end{figure*}

\FloatBarrier
\subsection{Operating thresholds and decision-curve analysis}
Thresholds fixed on development-set out-of-fold predictions transferred reasonably to the hold-out set (Table~\ref{tab:thresholds}). For stacking + RFECV, the 80\%, 85\% and 90\% sensitivity targets gave hold-out sensitivities of 0.815, 0.852 and 0.907 and specificities of 0.66, 0.63 and 0.56, flagging 51\%--60\% of patients for confirmatory testing. The 95\% target gave a hold-out sensitivity of 0.926 (0.848--0.983), below target. With only 54 positive patients, each missed case changes sensitivity by 1.9 percentage points. Logistic regression gave very similar operating points; at the 95\% target it reached 0.963 with specificity 0.53. NPV was 0.87--0.96 across all thresholds and both models.

In this internal decision-curve analysis (Fig.~\ref{fig:dca}), all three models had higher net benefit than treating all or no patients across threshold probabilities from 0.10 to 0.55; logistic regression fell to zero net benefit at 0.60. The three curves were nearly identical up to a threshold of 0.45 (differences in net benefit of at most 0.03). Above 0.50, the stack retained slightly more net benefit than logistic regression, but this region is based on few patients.

\begin{table*}[!htbp]
\centering
\caption{Sensitivity-targeted operating points. Thresholds were chosen on development-set out-of-fold predictions and applied unchanged to the hold-out set. Bootstrap 95\% CIs are shown in grey beneath. ``Flagged'' is the proportion of hold-out patients referred for confirmatory testing.}
\label{tab:thresholds}
\footnotesize
\setlength{\tabcolsep}{4pt}
\resizebox{\textwidth}{!}{%
\begin{tabular}{lccccccc}
\toprule
Model & Target & Threshold & Sensitivity & Specificity & PPV & NPV & Flagged \\
\midrule
Stacking + RFECV & 80\% & 0.433 & 0.815 & 0.660 & 0.564 & 0.868 & 50.6\% \\
 &  &  & {\scriptsize\color{gray}(0.709--0.915)} & {\scriptsize\color{gray}(0.559--0.747)} & {\scriptsize\color{gray}(0.453--0.670)} & {\scriptsize\color{gray}(0.790--0.940)} &  \\[2pt]
Stacking + RFECV & 85\% & 0.348 & 0.852 & 0.630 & 0.554 & 0.887 & 53.9\% \\
 &  &  & {\scriptsize\color{gray}(0.754--0.942)} & {\scriptsize\color{gray}(0.529--0.723)} & {\scriptsize\color{gray}(0.442--0.658)} & {\scriptsize\color{gray}(0.809--0.957)} &  \\[2pt]
Stacking + RFECV & 90\% & 0.276 & 0.907 & 0.560 & 0.527 & 0.918 & 60.4\% \\
 &  &  & {\scriptsize\color{gray}(0.821--0.981)} & {\scriptsize\color{gray}(0.459--0.653)} & {\scriptsize\color{gray}(0.420--0.622)} & {\scriptsize\color{gray}(0.839--0.983)} &  \\[2pt]
Stacking + RFECV & 95\% & 0.212 & 0.926 & 0.520 & 0.510 & 0.929 & 63.6\% \\
 &  &  & {\scriptsize\color{gray}(0.848--0.983)} & {\scriptsize\color{gray}(0.423--0.621)} & {\scriptsize\color{gray}(0.408--0.606)} & {\scriptsize\color{gray}(0.855--0.984)} &  \\[2pt]
Logistic regression & 80\% & 0.413 & 0.852 & 0.650 & 0.568 & 0.890 & 52.6\% \\
 &  &  & {\scriptsize\color{gray}(0.754--0.942)} & {\scriptsize\color{gray}(0.551--0.738)} & {\scriptsize\color{gray}(0.458--0.671)} & {\scriptsize\color{gray}(0.814--0.958)} &  \\[2pt]
Logistic regression & 85\% & 0.345 & 0.852 & 0.620 & 0.548 & 0.886 & 54.5\% \\
 &  &  & {\scriptsize\color{gray}(0.754--0.942)} & {\scriptsize\color{gray}(0.520--0.710)} & {\scriptsize\color{gray}(0.438--0.654)} & {\scriptsize\color{gray}(0.806--0.956)} &  \\[2pt]
Logistic regression & 90\% & 0.284 & 0.889 & 0.580 & 0.533 & 0.906 & 58.4\% \\
 &  &  & {\scriptsize\color{gray}(0.804--0.968)} & {\scriptsize\color{gray}(0.480--0.674)} & {\scriptsize\color{gray}(0.427--0.633)} & {\scriptsize\color{gray}(0.829--0.970)} &  \\[2pt]
Logistic regression & 95\% & 0.212 & 0.963 & 0.530 & 0.525 & 0.964 & 64.3\% \\
 &  &  & {\scriptsize\color{gray}(0.903--1.000)} & {\scriptsize\color{gray}(0.431--0.630)} & {\scriptsize\color{gray}(0.419--0.624)} & {\scriptsize\color{gray}(0.903--1.000)} &  \\[2pt]
\bottomrule

\end{tabular}}
\end{table*}

\begin{figure}[!htbp]
  \centering
  \includegraphics[width=0.6\linewidth]{fig6_dca.pdf}
  \caption{Decision curves on the hold-out set. Net benefit is shown against the threshold probability at which a patient would be referred for confirmatory testing.}
  \label{fig:dca}
\end{figure}

\FloatBarrier
\subsection{Sensitivity to missing-data handling}
The conclusions were robust to the imputation strategy (Table~\ref{tab:missing}). With iterative or $k$-NN imputation, ROC-AUC changed by at most $+0.005$ relative to median imputation for both logistic regression and RFECV + stacking (all $p \geq 0.45$). Adding missingness indicators did not help. It lowered the ROC-AUC of logistic regression slightly ($-0.010$, $p = 0.39$) and left RFECV + stacking unchanged ($-0.002$, $p = 0.50$), and RFECV retained on average only 1.2 of the five indicators. Under every strategy, stacking and logistic regression remained within 0.005 of each other.

\begin{table*}[!htbp]
\centering
\caption{Missing-data sensitivity analysis on the development set (10-fold cross-validation repeated twice; 20 folds). ROC-AUC values are means with corrected 95\% CIs. $\Delta$AUC is the paired difference from median imputation for the same model, with the corrected resampled $t$-test $p$-value.}
\label{tab:missing}
\resizebox{\textwidth}{!}{%
\begin{tabular}{llcccccl}
\toprule
Model & Imputation & Features & ROC-AUC (95\% CI) & MCC & Sens. & Brier & $\Delta$AUC ($p$) \\
\midrule
Logistic regression & Median & 16.0 & 0.846 (0.813--0.878) & 0.514 & 0.743 & 0.162 & reference \\
Logistic regression & Median + indicators & 21.0 & 0.836 (0.788--0.884) & 0.484 & 0.727 & 0.166 & -0.010 (0.387) \\
Logistic regression & Iterative (MICE-style) & 16.0 & 0.847 (0.814--0.879) & 0.517 & 0.753 & 0.162 & +0.001 (0.712) \\
Logistic regression & KNN ($k=10$) & 16.0 & 0.847 (0.813--0.880) & 0.505 & 0.743 & 0.162 & +0.001 (0.607) \\
\midrule
Stacking + RFECV & Median & 11.3 & 0.844 (0.804--0.883) & 0.498 & 0.752 & 0.164 & reference \\
Stacking + RFECV & Median + indicators & 13.3 & 0.841 (0.799--0.884) & 0.496 & 0.748 & 0.166 & -0.002 (0.503) \\
Stacking + RFECV & Iterative (MICE-style) & 14.0 & 0.848 (0.815--0.880) & 0.512 & 0.762 & 0.163 & +0.004 (0.594) \\
Stacking + RFECV & KNN ($k=10$) & 13.6 & 0.849 (0.814--0.884) & 0.501 & 0.759 & 0.163 & +0.005 (0.448) \\
\bottomrule

\end{tabular}}
\end{table*}


\FloatBarrier
\subsection{Explanation of predictions}
Glucose dominated the SHAP explanations of the stacking + RFECV pipeline (mean $|\mathrm{SHAP}| = 0.170$ on the probability scale), followed by BMI (0.074), age (0.055), the diabetes pedigree function (0.038) and pregnancies (0.028) (Fig.~\ref{fig:shap}). Insulin, skin-fold thickness and blood pressure contributed little (at most 0.015). Higher glucose, BMI, age and pedigree values consistently increased the predicted probability. These inputs contributed most to the model's predictions, in directions that agree with established risk factors; as model attributions, they do not measure causal effects. Because the outcome itself is largely defined by 2-h post-load glucose, the dominance of glucose in particular should not be interpreted as evidence of independent causal importance.

\begin{figure}[!htbp]
  \centering
  \includegraphics[width=0.8\linewidth]{fig7_shap.pdf}
  \caption{SHAP values of the hold-out predictions of the stacking + RFECV pipeline, expressed in the eight raw clinical inputs. Each point is one patient; colour shows the feature value.}
  \label{fig:shap}
\end{figure}
